% =============================================================================
% CHAPTER 6 — IMPLEMENTATION
% Thesis: PATHCAST
% Status: preliminary draft grounded in PathCast v0.1.0 harvest (Apr 2026)
% =============================================================================

\chapter{Implementation}
\label{ch:implementation}

This chapter documents the software instantiation of PATHCAST
(Definition~\ref{def:pipeline}). At the time of writing, Stage~1 (event log
construction) is implemented and exercised on the evaluation corpus described
in Chapter~\ref{ch:methodology}; Stages~2--4 and the ML component remain
specified but not yet integrated into the executable pipeline
(Section~\ref{sec:impl-roadmap}).

Section~\ref{sec:impl-architecture} presents the repository layout and
technology stack. Section~\ref{sec:impl-stage1} details the Stage~1 modules
and reports harvest-scale statistics from the April~2026 snapshot.
Section~\ref{sec:impl-artifacts} lists the interim data products.
Section~\ref{sec:impl-e1-pilot} reports a single-repository E1 feasibility
pilot. Section~\ref{sec:impl-confidential} discusses the handling of
confidential organisational data used for internal validation only.


% =============================================================================
% 6.1 ARCHITECTURE
% =============================================================================
\section{Software Architecture}
\label{sec:impl-architecture}

The implementation is organised as a Python package (\texttt{pathcast}, version
0.1.0) under a standard \texttt{src/} layout, with configuration in YAML
(\texttt{configs/repositories.yaml}, \texttt{configs/stage1\_github.yaml}) and
CLI scripts under \texttt{scripts/}. Table~\ref{tab:impl-modules} maps thesis
pipeline stages to code modules.

\begin{table}[htbp]
\centering
\caption{PATHCAST implementation modules (v0.1.0).}
\label{tab:impl-modules}
\small
\begin{tabular}{llll}
\toprule
\textbf{Stage} & \textbf{Package} & \textbf{Status} & \textbf{Entry point} \\
\midrule
$\mathcal{S}_1$ & \texttt{stage1\_eventlog} & Implemented & \texttt{run\_stage1\_github.py} \\
$\mathcal{S}_2$ & \texttt{stage2\_mining}    & Implemented (E1) & \texttt{run\_e1\_batch.py} \\
$\mathcal{S}_3$ & \texttt{stage3\_markov}    & Scaffold only & --- \\
$\mathcal{S}_4$ & \texttt{stage4\_montecarlo}& Scaffold only & --- \\
$\mathcal{S}_{\text{ML}}$ & \texttt{ml}        & Scaffold only & --- \\
Evaluation      & \texttt{evaluation}        & Pilot script  & \texttt{run\_e1\_pilot\_pm4py.py} \\
\bottomrule
\end{tabular}
\end{table}

Core dependencies include PM4Py~2.7 for event-log manipulation and conformance
checking, Pandas/Parquet for columnar storage, PyGithub and Requests for REST
collection, and Loguru for structured logging. Unit tests under
\texttt{tests/unit/} cover the XES builder activity mapping and edge cases
(eight tests at v0.1.0).


% =============================================================================
% 6.2 STAGE 1
% =============================================================================
\section{Stage 1: Event Log Construction}
\label{sec:impl-stage1}

Stage~1 implements the signature $\mathcal{S}_1 : \mathcal{R} \to \mathcal{L}$
(Equation~\ref{eq:stage1-sig}). The harvest proceeds in three steps:
\begin{enumerate}
    \item \textbf{Collection.} \texttt{collector.py} retrieves issues, timeline
    events, pull requests, and (optionally) GitHub Actions runs via the GitHub
    REST API, with pagination, rate-limit handling, and retry logic.
    \item \textbf{Normalisation.} Issue and PR histories are mapped to the
    PATHCAST activity alphabet (e.g.\ \texttt{ISSUE\_OPENED},
    \texttt{PR\_MERGED}; see \texttt{xes\_builder.py}).
    \item \textbf{Serialisation.} Events are written as a columnar Parquet file
    and, on demand, as an XES~IEEE~1849-2016 log via PM4Py.
\end{enumerate}

For Jira-backed repositories (Apache and Spring projects), \texttt{jira\_collector.py}
reads transition histories from the Montgomery et al.\ public Jira dataset and
produces a separate interim log (\texttt{jira\_event\_log.parquet}, 1.37~M events).

\subsection{Preliminary harvest statistics}
\label{sec:impl-stage1-stats}

Table~\ref{tab:impl-harvest} summarises the GitHub Stage-1 harvest executed in
April~2026 over 43 of the 48 evaluation repositories (five repositories
pending or empty at harvest time; see Appendix~\ref{app:repository-list}).

\begin{table}[htbp]
\centering
\caption{Stage-1 GitHub harvest snapshot (April 2026, $n = 43$ repositories).}
\label{tab:impl-harvest}
\begin{tabular}{lr}
\toprule
\textbf{Quantity} & \textbf{Value} \\
\midrule
Total events                         & 174{,}036 \\
Total process cases (issues + PRs)   & 24{,}066 \\
Distinct activity types (global)     & 18 \\
Mean events per repository         & 4{,}047 \\
Median events per repository       & 164 \\
Mean cases per repository          & 560 \\
Median cases per repository        & 82 \\
\midrule
Largest repository by cases          & \texttt{cli/cli} (9{,}548 cases) \\
Second largest                       & \texttt{fastapi/fastapi} (8{,}693) \\
\bottomrule
\end{tabular}
\end{table}

The highly skewed distribution (median $\ll$ mean) reflects the stratified
corpus: a few high-activity repositories (\texttt{cli/cli},
\texttt{fastapi/fastapi}) dominate event volume, while many medium-scale
projects contribute fewer than 200 cases each. This heterogeneity is intentional
(Section~\ref{sec:dataset-rationale}) and will be preserved in the full E1/E2
reporting by repository-level aggregation.


% =============================================================================
% 6.3 ARTIFACTS
% =============================================================================
\section{Interim Data Artifacts}
\label{sec:impl-artifacts}

Table~\ref{tab:impl-artifacts} lists the reproducible artefacts produced by
the current implementation. Raw API downloads reside under
\texttt{data/raw/github/}; processed Stage-1 outputs under
\texttt{data/interim/stage1/}.

\begin{table}[htbp]
\centering
\caption{Interim artefacts at v0.1.0.}
\label{tab:impl-artifacts}
\small
\begin{tabular}{llr}
\toprule
\textbf{File} & \textbf{Description} & \textbf{Size (approx.)} \\
\midrule
\texttt{github\_event\_log.parquet} & GitHub Stage-1 log (43 repos) & 2.9~MB \\
\texttt{jira\_event\_log.parquet}   & Montgomery Jira transitions  & 35~MB \\
\texttt{jira\_event\_log.xes}       & XES export (Jira)            & 645~MB \\
\texttt{dataset\_summary\_stats.json}& Auto-computed table inputs  & $<$1~KB \\
\bottomrule
\end{tabular}
\end{table}

Aggregate dataset statistics derived from these artefacts appear in
Table~\ref{tab:dataset-summary} (Chapter~\ref{ch:methodology}).


% =============================================================================
% 6.4 E1 PILOT
% =============================================================================
\section{Preliminary E1 Feasibility Pilot}
\label{sec:impl-e1-pilot}

Although Stage~2 is not yet packaged as a library module, a standalone pilot
script (\texttt{scripts/run\_e1\_pilot\_pm4py.py}) exercises the E1 protocol of
Section~\ref{sec:e1-design} on a single public repository using PM4Py directly.
The pilot filters the 800 most recent cases from \texttt{fastapi/fastapi},
discovers a Petri net with the Inductive Miner at
$\theta_{\text{noise}} = 0.2$, and computes alignment-based fitness and
precision plus token-based-replay generalization.

% Auto-generated companion to E1 pilot — update after scripts/run_e1_pilot_pm4py.py
\begin{table}[htbp]
\centering
\caption{Preliminary E1 pilot on \texttt{fastapi/fastapi} (800 most recent cases;
Inductive Miner, $\theta_{\text{noise}} = 0.2$). Not the full 48-repository E1.}
\label{tab:e1-pilot}
\small
\begin{tabular}{lrr}
\toprule
\textbf{Metric} & \textbf{Value} & \textbf{E1 threshold} \\
\midrule
Cases analysed            & 800   & $\geq 50$ \\
Events                    & 6{,}251 & --- \\
Distinct activities       & 14    & --- \\
Activity coverage         & 1.00  & $\geq 0.50$ \\
Petri net places          & 41    & --- \\
Petri net transitions     & 52    & --- \\
Fitness (alignments)      & 0.911 & $\geq 0.70$ \\
Precision (alignments)    & 0.817 & $\geq 0.50$ \\
Generalization (TBR)      & 0.728 & --- \\
\bottomrule
\multicolumn{3}{l}{\parbox{11cm}{\footnotesize Single-repository feasibility
demonstration only; full E1 aggregates over all repositories satisfying
Section~\ref{sec:e1-design}.}}
\end{tabular}
\end{table}


\subsection{E1 batch on ten representative repositories}
\label{sec:impl-e1-batch}

Table~\ref{tab:e1-representative} reports the preliminary E1 batch executed with
\texttt{scripts/run\_e1\_batch.py} on the ten stratified repositories defined in
Section~\ref{sec:e1-design} (sensitivity subset). All ten repositories satisfy
the preview acceptance criteria of Equation~\eqref{eq:acceptance}.

% Auto-generated by scripts/run_e1_batch.py
\begin{table}[htbp]
\centering
\caption{E1 preliminary batch on ten representative repositories (IMf, $\theta_{\text{noise}} = 0.2$, up to 800 cases each).}
\label{tab:e1-representative}
\small
\begin{tabular}{llrrrr}
\toprule
\textbf{Repository} & \textbf{Source} & \textbf{Cases} & \textbf{Fitness} & \textbf{Precision} & \textbf{Pass} \\
\midrule
cli/cli & GitHub & 800 & 0.886 & 0.678 & Yes \\
fastapi/fastapi & GitHub & 800 & 0.911 & 0.786 & Yes \\
django/django & GitHub & 62 & 1.000 & 1.000 & Yes \\
facebook/react & GitHub & 84 & 1.000 & 1.000 & Yes \\
golang/go & GitHub & 187 & 0.990 & 0.764 & Yes \\
rust-lang/rust & GitHub & 410 & 0.989 & 0.998 & Yes \\
microsoft/vscode & GitHub & 758 & 0.993 & 0.994 & Yes \\
scikit-learn/scikit-learn & GitHub & 82 & 0.997 & 1.000 & Yes \\
apache/kafka & Jira & 800 & 0.911 & 0.583 & Yes \\
spring-projects/spring-framework & Jira & 800 & 0.916 & 0.620 & Yes \\
\bottomrule
\multicolumn{6}{l}{\parbox{14cm}{\footnotesize Pass = preview acceptance criteria of Eq.~\eqref{eq:acceptance} (fitness, precision, coverage, cases).}}
\end{tabular}
\end{table}


All four preview acceptance checks of Equation~\eqref{eq:acceptance} are
satisfied for this repository (fitness $= 0.911$, precision $= 0.817$,
coverage $= 1.00$, cases $= 800$). This result supports the \emph{feasibility}
of the E1 design on a high-activity GitHub Issues workflow but must not be
extrapolated to the full 48-repository evaluation: a single pilot cannot
establish external validity, and several repositories in the corpus use Jira
workflows with different activity alphabets.


% =============================================================================
% 6.5 CONFIDENTIAL DATA
% =============================================================================
\section{Confidential Organisational Data}
\label{sec:impl-confidential}

During internal development, event logs extracted from a confidential
organisational environment (Jira and Bitbucket projects) were used
\emph{off-line} to stress-test Stage~1 activity mapping and PM4Py
integration. These data are subject to organisational terms of use and are
not part of the public evaluation corpus (Chapter~\ref{ch:methodology}).

\paragraph{Anonymisation assessment.}
A field-level review of the raw exports identifies the following
re-identification risks if published without transformation:
\begin{itemize}
    \item \textbf{Direct identifiers:} employee names in the \texttt{Author}
    column; Atlassian URLs containing the organisation domain; issue titles
    describing internal deliverables; and project keys.
    \item \textbf{Quasi-identifiers:} exact timestamps at second resolution;
    rare activity sequences combined with status labels customised to the
    organisation's workflow.
    \item \textbf{Lower risk (aggregates):} DFG edge counts, bottleneck tables
    by status name, fitness/precision on pseudonymised logs, and variant
    \emph{structure} without titles or authors.
\end{itemize}

\paragraph{Thesis reporting policy.}
Unless explicit written authorisation is obtained from the data controller,
\emph{no record from the confidential corpus appears in this thesis or in the
public replication package}. Permitted uses without authorisation are limited to:
\begin{enumerate}
    \item methodological prose describing how confidential logs \emph{would} be
    handled (pseudonymisation, URL removal, date shifting);
    \item fully synthetic or publicly available substitute logs for figures; and
    \item aggregate-only statistics that cannot be linked to the organisation
    (e.g.\ ``a private Jira project with 987 cases and 36 activities yielded
    fitness $\geq 0.7$''), provided no workflow/status names unique to that
    organisation are quoted.
\end{enumerate}

Even when individual status labels appear generic, the combination of labels,
volumes, and timing could still fingerprint the organisation; therefore the
conservative default is non-disclosure until a formal anonymisation protocol
and legal review are completed.


% =============================================================================
% 6.6 ROADMAP
% =============================================================================
\section{Implementation Roadmap}
\label{sec:impl-roadmap}

The remaining milestones follow the PathCast development roadmap (v0.2--v0.6):
Stage~2 process discovery and conformance as a library module; Stage~3 Markov
estimation with Laplace smoothing for sparse transitions; Stage~4 Monte Carlo
simulation with percentile extraction; orchestrating CLI; and the full E1/E2
evaluation harness. Upon completion, this chapter will be expanded with
end-to-end runtime metrics, replication-package DOI, and PostgreSQL persistence
(Section~\ref{sec:data-collection}).
